% !PRINTED 129
\weberpage{129}
%<<CONT

\begin{equation}
 f(\alpha+h)=f(\alpha)+h f'(\alpha)+\frac{h^2}{1\cdot2}f''(\alpha)+\cdots
\tag{6}
\end{equation}
\begin{equation}
 h=-\delta\frac{f(\alpha)}{f'(\alpha)},
\tag{7}
\end{equation}
and obtain
\begin{align}
 f(\alpha+h)={}&(1-\delta)f(\alpha)
 +\delta^2\left\{\frac{f(\alpha)^2 f''(\alpha)}{1\cdot2\,f'(\alpha)^2}
 -\delta\frac{f(\alpha)^3 f'''(\alpha)}{1\cdot2\cdot3\,f'(\alpha)^3}+\cdots\right\}.
\tag{8}
\end{align}

Now by (2)
\[
 f'(\alpha)=n(\alpha-\beta_1)(\alpha-\beta_2)\cdots(\alpha-\beta_{n-1});
\]
therefore, if the absolute values of $(\alpha-\beta_1),(\alpha-\beta_2),
\ldots,(\alpha-\beta_{n-1})$ are denoted by $r_1,r_2,
\ldots,r_{n-1}$,
\[
 |f'(\alpha)|=nr_1r_2\cdots r_{n-1}.
\]
Since $\alpha$ lies outside the circle $\varrho_1$ drawn about $\beta_1$, we have $r_1>\varrho_1,r_2>\varrho_2,\ldots$, and consequently
\[
 |f'(\alpha)|>n\varrho_1\varrho_2\cdots\varrho_{n-1}.
\]
Thus $|f'(\alpha)|$ is greater than a positive number $k$ independent of the position of $\alpha$ in $G$. It follows that one can choose a sufficiently large positive number $Q$, likewise independent of the position of $\alpha$ and of the proper fraction $\delta$, so that
\begin{equation}
 \left|\frac{f(\alpha)^2f''(\alpha)}{1\cdot2\,f'(\alpha)^2}
 -\delta\frac{f(\alpha)^3f'''(\alpha)}{1\cdot2\cdot3\,f'(\alpha)^3}+\cdots\right|<Q .
\tag{9}
\end{equation}

One can obtain $Q$ by replacing in the denominators on the left of (9) $f'(\alpha)$ by $k$, replacing $\delta$ by $1$, replacing all the other terms by their absolute values, replacing finally the absolute value of $\alpha$ by the greater value $R$, and then making the number so obtained arbitrarily larger. From (8) one then obtains
\begin{equation}
 |f(\alpha+h)|<(1-\delta)|f(\alpha)|+\delta^2Q.
\tag{10}
\end{equation}
Assume now $Q$ so chosen that, for every $\alpha$ in $G$,
\[
 2Q>|f(\alpha)|=a
\]
which is plainly compatible with the above determination. We may then put
\[
 \delta=\frac{a}{2Q},\qquad
 h=-\frac{a f(\alpha)}{2Q f'(\alpha)}
\]
%CONT>>
